\begin{remark}[Influence of the spectrally separated state]
The first two eigenvectors alone imply only
$\kappa_*(A_\varepsilon)\geq4\varepsilon^{-2}[1+o(1)]$.
The sharp constant eight appears only after the third orthogonality constraint
is incorporated.  This provides a concrete metric-optimization analogue of the
broader observation that states away from a coalescing pair can still alter
non-Hermitian eigenstate geometry near an exceptional point
\cite{Schomerus2024}.  The quantities are different: the cited work concerns
eigenvalue sensitivity and phase rigidity, whereas the present statement is
about the minimum spectral condition number among positive quasi-Hermitian
metrics.
\end{remark}

\subsection{A uniform perturbation bound for normalized solutions}
The metric certificate collapses in this family, but a conventional
resolvent estimate gives a uniform bound on inversion and on the
normalized solution state for all sufficiently small Euclidean errors.

\begin{proposition}[Euclidean stability independent of metric conditioning]
\label{prop:uniform_state}
Let $A=A_\varepsilon$, $0<\varepsilon\leq1/2$, and set
$M=\sqrt{15}/2$. For $\delta=\|\Delta\|_2<M^{-1}$, both $A$ and
$A+\Delta$ are invertible. For every Euclidean-unit vector $b$, put
$y=A^{-1}b$ and $z=(A+\Delta)^{-1}b$. Then
\begin{equation}\label{eq:uniform_state}
 \frac{\|z-y\|_2}{\|y\|_2}
 \leq\frac{M\delta}{1-M\delta},\qquad
 \left\|\frac{z}{\|z\|_2}-\frac{y}{\|y\|_2}\right\|_2
 \leq\frac{2M\delta}{1-M\delta}.
\end{equation}
The bounds do not depend on $\varepsilon$ or on an admissible metric.
\end{proposition}
\begin{proof}
Theorem~\ref{thm:counterexample} gives $\|A^{-1}\|_2\leq M$.
For $M\delta<1$, the Neumann series yields
$\|(A+\Delta)^{-1}\|_2\leq M/(1-M\delta)$.
The identity $z-y=-(A+\Delta)^{-1}\Delta y$ gives the first
inequality. For nonzero vectors $u,v$, adding and subtracting
$v/\|u\|_2$ gives
$\|u/\|u\|_2-v/\|v\|_2\|_2
\leq2\|u-v\|_2/\|v\|_2$; taking $u=z$, $v=y$
gives the second. The bound also controls the pure-state trace distance
(up to its trivial upper bound of one).
\end{proof}
